\documentclass[a4paper,12pt]{article}
\usepackage[utf8]{inputenc}
\usepackage[T1]{fontenc}
\usepackage[ngerman]{babel}
\usepackage{geometry}
\geometry{margin=2.5cm}
\usepackage{hyperref}
\usepackage{parskip}
\setlength{\parskip}{6pt}

\begin{document}

\section*{Learning to Simulate Complex Physics with Graph Networks}

\textbf{Autoren:} Alvaro Sanchez-Gonzalez, Jonathan Godwin, Tobias Pfaff, Rex Ying,
Jure Leskovec, Peter W.\ Battaglia \\
\textbf{Venue:} International Conference on Machine Learning (ICML), PMLR 119, 2020 \\
\textbf{Quelle:} \href{https://arxiv.org/abs/2002.09405}{arXiv:2002.09405}

% ------------------------------------------------------------------
\subsection*{1.\quad Kernidee und Motivation}
% ------------------------------------------------------------------

Traditionelle numerische Simulatoren physikalischer Systeme sind rechenintensiv,
schwer zu generalisieren und erfordern erheblichen Implementierungsaufwand. Die Arbeit
von Sanchez-Gonzalez et al.\ (2020) verfolgt daher den Ansatz, einen Simulator
\emph{direkt aus Daten zu erlernen}, anstatt explizite physikalische Gleichungen
numerisch zu lösen. Das vorgestellte Framework -- \emph{Graph Network-based Simulators}
(GNS) -- repräsentiert den Zustand eines physikalischen Systems durch Partikel als
Knoten eines Graphen und berechnet die zeitliche Dynamik mittels \emph{gelerntem
Message-Passing}.

Die zentrale Motivation liegt dabei in der Verallgemeinerungsfähigkeit: Ein einmal
trainiertes GNS-Modell soll auf Konfigurationen generalisieren, die während des
Trainings nicht gesehen wurden -- insbesondere auf größere Systeme mit einer deutlich
höheren Partikelanzahl sowie auf längere Zeithorizonte. Diese Eigenschaft rückt die
Arbeit in direkten Bezug zur Fragestellung der vorliegenden Masterarbeit.

% ------------------------------------------------------------------
\subsection*{2.\quad Architektur: Encode-Process-Decode}
% ------------------------------------------------------------------

Das GNS-Framework folgt einem dreistufigen \emph{Encode-Process-Decode}-Schema:

\begin{description}
  \item[Encoder] Der Encoder bildet den partikelbasierten Zustandsvektor $X$ auf
    einen latenten Graphen $G^0$ ab. Knotenmerkmale kodieren Positionen, eine
    Sequenz vergangener Geschwindigkeiten ($C = 5$ Zeitschritte) sowie
    materialspezifische Eigenschaften (Wasser, Sand, verformbares Material,
    Randbedingung). Kantenmerkmale kodieren relative Verschiebungsvektoren zwischen
    benachbarten Partikeln innerhalb eines Verbindungsradius $R$ -- eine Designwahl,
    die Translationsinvarianz erzwingt und sich als wesentlich für die
    Generalisierungsfähigkeit erweist.

  \item[Processor] Der Processor führt $M$ Schritte von gelernetem Message-Passing
    durch, wobei jeder Schritt durch ein eigenes Graph-Netzwerk mit ungeteilten
    Parametern realisiert wird. Residualverbindungen zwischen Ein- und Ausgabe jedes
    Message-Passing-Schritts stabilisieren das Training. Die Anzahl der Schritte $M$
    bestimmt maßgeblich die Modellkapazität, da größere $M$ die Kommunikation über
    größere räumliche Reichweiten ermöglichen.

  \item[Decoder] Der Decoder extrahiert aus den finalen Knotenrepräsentationen
    $v_i^M$ die vorhergesagten Beschleunigungen $\ddot{p}_i$, aus denen mittels
    semi-implizitem Euler-Integrator die nächsten Positionen und Geschwindigkeiten
    berechnet werden.
\end{description}

Das Training erfolgt auf einem \emph{Ein-Schritt-Verlust} (L2-Norm der
Beschleunigungen), ergänzt durch eine entscheidende Regularisierungsmaßnahme:
Trainingseingaben werden mit \emph{Random-Walk-Rauschen} korrupt, um die
Verteilungsverschiebung zwischen dem Training (saubere Eingaben) und dem
Rollout-Betrieb (akkumulierte Modellfehler als Eingaben) zu überbrücken.

% ------------------------------------------------------------------
\subsection*{3.\quad Zentrale Ergebnisse}
% ------------------------------------------------------------------

\paragraph{Generalisierung über Partikelzahlen.}
Ein auf dem WATERRAMPS-Datensatz (ca.\ 2500 Partikel pro Trajektorie) trainiertes
Modell generalisiert erfolgreich auf Szenen mit bis zu 85\,000 Partikeln -- eine
Steigerung um den Faktor 34 gegenüber dem Training. Rollouts über 5000 Zeitschritte
(8-fach länger als die Trainingstrajektorien) bleiben physikalisch plausibel. Dieses
Ergebnis demonstriert, dass der Einsatz \emph{relativer Positionen} als
Kantenmerkmal eine starke Transferfähigkeit auf nie gesehene Systemgrößen ermöglicht.

\paragraph{Materialübergreifendes Lernen.}
Ein einzelnes GNS-Modell kann simultan Wasser, Sand und viskoplastisches Material
(,,Goop``) simulieren, einschließlich ihrer Wechselwirkungen untereinander. Dabei
wird das Materialmerkmal lediglich als Eingangsfeature übergeben -- keine
architektonischen Anpassungen sind erforderlich.

\paragraph{Einfluss architektonischer Entscheidungen.}
Systematische Ablationen zeigen, dass folgende Faktoren die Modellperformance
signifikant beeinflussen: (i) die Anzahl der Message-Passing-Schritte $M$ (mehr
Schritte verbessern Langzeitrollouts durch längere Reichweite), (ii) ungeteilte
Gewichte im Processor (im Vergleich zu geteilten Gewichten), (iii) der
Verbindungsradius $R$ sowie (iv) die Verwendung relativer statt absoluter
Positionen. Hyperparameter wie MLP-Tiefe, Latentgröße und das Vorhandensein von
Selbstkanten haben hingegen nur geringen Einfluss.

\paragraph{Vergleich mit Baseline-Modellen.}
GNS übertrifft Continuous Convolution (CConv, Ummenhofer et al.\ 2020) und DPI
(Li et al.\ 2018) auf allen getesteten Datensätzen, obwohl CConv speziell für
Fluidsimulationen entwickelt wurde. GNS benötigt dabei keine materialspezifischen
Anpassungen.

% ------------------------------------------------------------------
\subsection*{4.\quad Kritische Einordnung}
% ------------------------------------------------------------------

Ein konzeptionell wichtiger Aspekt des GNS-Ansatzes ist die \emph{Markov-Annahme}:
Das Modell wird auf dem Ein-Schritt-Verlust optimiert, ohne explizit die
langfristige Trajektorie zu berücksichtigen. Die Autoren argumentieren, dass
diese Vereinfachung die Generalisierbarkeit gegenüber einer Optimierung auf
Rollout-Verlust sogar verbessern kann, da Rollout-Training das Modell dazu
verleiten kann, Bias-Korrekturen in den Voraussagen zu kodieren, die auf
spezifische Anfangsbedingungen zugeschnitten sind.

Darüber hinaus operiert GNS im \emph{zeitabhängigen} Regime (Navier-Stokes bei
moderaten Reynoldszahlen, sowie MPM-Simulationen), während die vorliegende
Masterarbeit \emph{stationäre} Stokes-Strömung ($Re \ll 1$) betrachtet. Dieser
Unterschied im physikalischen Regime hat fundamentale Konsequenzen für die
Wahl des Ausgabe-Targets, der Netzwerkarchitektur und des Trainingsverfahrens.

% ------------------------------------------------------------------
\subsection*{5.\quad Bezug zur Masterarbeit}
% ------------------------------------------------------------------

Das GNS-Paper berührt mehrere Kernfragen der vorliegenden Masterarbeit, obwohl
beide Arbeiten grundlegend unterschiedliche Paradigmen verfolgen:

\paragraph{Partikelbasiert vs.\ gitterbasiert.}
GNS kodiert physikalische Zustände partikelbasiert als Knoten eines variabel
aufgebauten Graphen. Dieser Ansatz ist natürlicherweise invariant gegenüber der
Anzahl und Anordnung von Kristallen -- eine Eigenschaft, die in der Masterarbeit
durch gitterbasierte Architekturen (U-Net, FNO) erreicht werden muss, indem
Kristallpositionen als Eingangsfeature auf ein festes $256 \times 256$-Gitter
kodiert werden. GNS bietet hierfür eine konzeptionell elegantere Lösung, allerdings
um den Preis eines erhöhten Implementierungsaufwands und der Schwierigkeit, auf
Gittergrößen direkt zu interpolieren.

\paragraph{Generalisierung über Konfigurationsanzahl.}
Die von GNS demonstrierte Generalisierung auf eine um den Faktor 34 höhere
Partikelzahl ist direkt relevant für die Hauptforschungsfrage der Masterarbeit:
Können Surrogat-Modelle, die auf Konfigurationen mit bis zu $N$ Kristallen
trainiert wurden, auf Konfigurationen mit $N + m$ Kristallen übertragen werden?
GNS löst dieses Problem durch die Wahl \emph{relativer Positionsrepräsentationen},
die invariant gegenüber der absoluten räumlichen Anordnung sind. In der Masterarbeit
entspricht dies der Wahl physikalisch motivierter Normalisierungsstrategien sowie
der Verwendung der Stromfunktion $\psi$ als Ausgabe-Target, da diese durch
Konstruktion Inkompressibilität garantiert.

\paragraph{Rauschen zur Fehlerakkumulierungskompensation.}
Die im GNS eingesetzte Strategie der Trainingsrauschkorruption (Random Walk) ist
konzeptionell verwandt mit der Datenaugmentierungsstrategie der Masterarbeit, mit der
die geometrische Variabilität des Trainingsdatensatzes von ca.\ 300 auf bis zu 1800
Samples erweitert werden soll. Beide Ansätze verfolgen das Ziel, die Verallgemeinerung
auf Verteilungen außerhalb der Trainingsdaten zu verbessern.

\paragraph{Physikalische Konsistenz vs.\ rein datengetriebene Vorhersage.}
GNS lernt keine physikalischen Erhaltungsgrößen explizit -- weder Divergenzfreiheit
noch Massenerhaltung sind architektonisch erzwungen. In der Masterarbeit wird
hingegen die Stromfunktionsformulierung gewählt, die Divergenzfreiheit ($\nabla
\cdot \mathbf{v} = 0$) durch mathematische Konstruktion auf Maschinengenauigkeit
($\approx 10^{-6}$) garantiert. Dies stellt einen wesentlichen Vorteil gegenüber
dem GNS-Ansatz für das Stokes-Strömungsregime dar, in dem Inkompressibilität eine
fundamentale physikalische Randbedingung ist.

\paragraph{Spektrale vs.\ räumliche Verarbeitung.}
Die im GNS identifizierte Schlüsselrolle der \emph{Reichweite} der
Partikel-Partikel-Interaktion -- quantifiziert durch den Verbindungsradius $R$ und
die Anzahl der Message-Passing-Schritte $M$ -- korrespondiert direkt mit einem
zentralen Vergleichsaspekt der Masterarbeit: Während der U-Net-Encoder lokale
räumliche Merkmale hierarchisch extrahiert, erlaubt das FNO in einem einzigen
Fourier-Layer den Zugriff auf globale Korrelationen über die gesamte Domäne.
Für elliptische Probleme wie Stokes-Strömung, bei denen lokale Störungen global
propagieren, ist diese globale Reichweite des FNO strukturell vorteilhaft --
ein Argument, das durch die GNS-Ablationsstudie zu $M$ und $R$ untermauert wird.

\paragraph{Mögliche Erweiterung: Graph-basierte Surrogat-Modelle.}
GNS eröffnet eine perspektivisch interessante Erweiterungsrichtung für zukünftige
Arbeiten: Ein partikelbasiertes Graph-Netzwerk, in dem Kristalle als Knoten eines
variablen Graphen kodiert werden, wäre von Natur aus in der Lage, auf beliebige
Kristallzahlen zu generalisieren -- ohne die Einschränkung einer festen Gittergröße.
Der vorliegende gitterbasierte Ansatz (U-Net / FNO) ist jedoch für die gestellte
Forschungsfrage geeigneter, da er eine direkte Vergleichbarkeit mit LaMEM-Referenzdaten
auf einem regulären Gitter ermöglicht und keine Interpolation zwischen Partikel-
und Gitterdarstellungen erfordert.

\end{document}